\chapter{Haemophilus Influenzae Type B (Hib)}
\index{Haemophilus Influenzae Type B (Hib)}

\section*{Definition and Overview}
Haemophilus influenzae type b (Hib) is a bacterium that can cause serious invasive infections, particularly in young children, including meningitis, pneumonia, epiglottitis, sepsis, and infections of the joints and bone. Before vaccination, Hib was the most common cause of bacterial meningitis in children. The Hib vaccine, which is part of the routine childhood immunisation schedule in many countries, has made Hib disease rare. Hib is spread through respiratory droplets and is only a risk to unvaccinated people, especially children under five, who are at highest risk of severe disease. Prompt treatment with antibiotics is essential for any invasive Hib infection.

\section*{Epidemiology}
Before the Hib vaccine was introduced in many countries in 1993, Hib caused around 500--700 cases of invasive disease per year, mostly in children under five, and was a leading cause of childhood meningitis and a common cause of epiglottitis. Since vaccination, the number of cases has fallen by more than 90\%. Current cases occur mainly in unvaccinated children, in children too young to have completed their vaccine course, and occasionally in older people with weakened immunity. Hib disease remains a major problem in countries without vaccination programs.

\section*{Aetiology and Pathophysiology}
Hib causes disease by invading the bloodstream from the upper respiratory tract.

\begin{itemize}
    \item \textbf{The bacterium:} Haemophilus influenzae type b has a polysaccharide capsule that helps it evade the immune system
    \item \textbf{Colonisation:} the bacterium lives harmlessly in the nose and throat of many people, mostly without causing disease
    \item \textbf{Invasion:} in susceptible people, particularly unvaccinated children under five, the bacterium enters the bloodstream
    \item \textbf{Dissemination:} from the blood, it spreads to the meninges (meningitis), lungs (pneumonia), epiglottis (epiglottitis), joints (septic arthritis), bone, or tissues (cellulitis)
    \item \textbf{Transmission:} through respiratory droplets from coughing, sneezing, or close contact
    \item \textbf{Immunity:} the vaccine generates protective antibodies against the capsule
\end{itemize}

\section*{Clinical Features}
The features depend on the site of infection and develop acutely.

\begin{itemize}
    \item Meningitis: fever, headache, stiff neck, vomiting, lethargy, irritability, and a bulging fontanelle in infants; can progress to seizures and coma
    \item Epiglottitis: sudden high fever, severe sore throat, drooling, difficulty swallowing, and breathing difficulty, which can obstruct the airway rapidly
    \item Pneumonia: fever, cough, and breathing difficulty
    \item Sepsis: fever, poor feeding, lethargy, and features of shock
    \item Septic arthritis: a hot, swollen, painful joint
    \item Cellulitis: a red, swollen, tender area of skin, often on the face or cheek
    \item Bone infection (osteomyelitis): pain, swelling, and fever
\end{itemize}

\section*{Differential Diagnosis}
Other causes of these serious infections should be considered.

\begin{itemize}
    \item Meningitis from other bacteria, including meningococcus, pneumococcus, and group B streptococcus
    \item Meningitis and encephalitis from viruses
    \item Epiglottitis from other bacteria and other causes of airway obstruction
    \item Pneumonia from viruses and other bacteria
    \item Sepsis from other organisms
    \item Septic arthritis and osteomyelitis from other bacteria
    \item Croup and other causes of stridor, which differ from epiglottitis
\end{itemize}

\section*{When to Seek Medical Care}
Seek urgent medical care for any child with high fever, lethargy, breathing difficulty, or signs of serious infection. A child who has been in contact with someone with Hib disease should be assessed, and preventive antibiotics may be recommended.

\textbf{Contact your local emergency services for:}
\begin{itemize}
    \item Difficulty breathing, drooling, or stridor, which may indicate epiglottitis
    \item Stiff neck, high fever, vomiting, and drowsiness, which may indicate meningitis
    \item A child who is difficult to wake or having a seizure
    \item Signs of shock, including pale, cold skin and rapid breathing
\end{itemize}

\section*{Diagnosis and Investigations}
Invasive Hib disease is diagnosed by cultures and rapid tests.

\begin{itemize}
    \item \textbf{Blood cultures:} detect the bacterium in the bloodstream
    \item \textbf{Lumbar puncture:} cerebrospinal fluid is examined for meningitis
    \item \textbf{Other cultures:} joint fluid, pleural fluid, or pus as appropriate
    \item \textbf{Antigen tests:} rapid detection of Hib in body fluids
    \item \textbf{Imaging:} chest X-ray for pneumonia, and imaging for joint or bone infection
    \item \textbf{Notification:} invasive Hib disease is a notifiable condition in many countries
    \item \textbf{Contact tracing:} household contacts may need antibiotics and vaccination assessment
\end{itemize}

\section*{Management}
Invasive Hib disease requires urgent hospital treatment.

\begin{itemize}
    \item \textbf{Intravenous antibiotics:} prompt treatment with antibiotics, usually third-generation cephalosporins
    \item \textbf{Intensive care:} for meningitis, epiglottitis, and severe sepsis
    \item \textbf{Airway management:} epiglottitis may require intubation to secure the airway
    \item \textbf{Supportive care:} fluids, oxygen, and management of complications
    \item \textbf{Surgery:} drainage of infected joints or abscesses
    \item \textbf{Prevention in contacts:} household contacts receive antibiotics to prevent spread
    \item \textbf{Vaccination:} ensuring the child's immunisations are up to date after recovery
    \item \textbf{Follow-up:} hearing and developmental assessment after meningitis
\end{itemize}

\section*{Complications}
\begin{itemize}
    \item Meningitis complications: hearing loss, seizures, learning difficulties, and neurological damage
    \item Airway obstruction from epiglottitis, which can be fatal without urgent care
    \item Sepsis and septic shock
    \item Septic arthritis and osteomyelitis with permanent joint or bone damage
    \item Pericarditis and other complications
    \item Death, which occurs in a small proportion of cases despite treatment
    \item Long-term disability in survivors of meningitis
\end{itemize}

\section*{Prognosis and Outcomes}
With prompt treatment, most children with invasive Hib disease recover, but it remains a serious illness. Meningitis has a mortality of around 5\% with treatment and causes lasting disability, particularly hearing loss, in 10--30\% of survivors. Epiglottitis is life-threatening if the airway is not secured quickly. Because the Hib vaccine is so effective, the best "treatment" is prevention: fully vaccinated children are protected against this disease, and the prognosis for vaccinated populations has transformed.

\section*{Prevention and Self-Care}
\begin{itemize}
    \item Ensure children receive the Hib vaccine on schedule: at 2, 4, and 6 months, with a booster at 18 months in many countries
    \item Catch up on any missed doses as soon as possible
    \item Check that your family's vaccinations are up to date, including for travel
    \item Practise good hand hygiene and cough etiquette
    \item Keep infants away from people with respiratory infections
    \item Seek urgent care for any child with signs of serious infection
    \item If a household contact develops Hib disease, follow public health advice about antibiotics and vaccination
    \item Breastfeed if possible, which supports infant immunity
\end{itemize}

\section*{Sources and Further Reading}
The following reputable sources provide further information for healthcare students and patients seeking reliable, current advice about Hib.

\begin{itemize}
    \item Australian Government Department of Health and Aged Care. \textit{Hib vaccine and immunisation schedule.} https://www.health.gov.au/
    \item National Centre for Immunisation Research and Surveillance. \textit{Hib vaccine fact sheet.} https://www.ncirs.org.au/
    \item Healthdirect Australia. \textit{Hib (Haemophilus influenzae type b).} https://www.healthdirect.gov.au/
    \item NSW Health. \textit{Hib disease fact sheet.} https://www.health.nsw.gov.au/
    \item Centers for Disease Control and Prevention. \textit{Hib disease.} https://www.cdc.gov/
    \item World Health Organization. \textit{Haemophilus influenzae type b.} https://www.who.int/
\end{itemize}
